\honest{Fake Morality}{Eliezer Yudkowsky}{2007}

Religious fundamentalists hold that there can be no morality without a Judge who rewards and punishes. If we did not fear hell and yearn for heaven, ``what would stop people from murdering each other left and right?'' \nb{The view has a long history. William Paley held that we need to believe in God as the final granter of happiness if we are to keep the motivation to do what we know we ought to do.}

Suppose Omega threatens to kill you if you ever use a bathroom between 7 and 10 in the morning. If Omega withdrew the threat, you would not panic. You would be relieved.

So a believer who is afraid of God no longer punishing murder has shown that they object to murder for reasons of their own. Otherwise the prospect would be no more horrifying than God not punishing sneezing. \nb{Plato's \textsc{Euthyphro} asks an older version of the question: whether the gods love the holy because it is holy, or it is holy because they love it. The opening question was about what would stop ``people''; the answer is about the believer who is afraid. It shows that this believer has a motive of their own. Whether people in general would keep one, which was Paley's worry, the post does not take up.}

To any religious readers of Overcoming Bias who are left: if you lose your faith, you will not lose your moral direction. The fear of losing a moral compass is itself a moral compass. ``I suspect you are steering by that compass, and that you always have been.''

You do not hear fundamentalists ask what would stop people from eating pork if there were no hell. Yet on their assumption, ``that we have no moral compass but divine reward and retribution,'' that argument should sound just as strong. \nb{The Talmud draws the same line. It separates laws that ``it would have been logical'' to write even if they had not been written, among them the bans on bloodshed and theft, from statutes whose reasons ``are not known,'' among them the ban on pork. So the assumption the post ascribes to fundamentalists is not a religious view as such.}

Even a threat of hellfire works only because people already dislike hellfire.

Then two philosophers of my own invention. The first says: be selfish, because people who try to improve society meddle and make everyone unhappy, and the market pays most for the work that benefits society most. The second says: be altruistic, because cooperation pays and people who contribute are happier. The first justifies selfishness by its benefit to everyone; the second justifies altruism by its benefit to himself.

Which of them is the real altruist? ``Whichever one actually holds open doors for little old ladies.''
